\documentclass[10pt,a4paper]{extarticle}
\usepackage[T1]{fontenc}
\usepackage[spanish]{babel}
\usepackage{graphicx}
\usepackage{booktabs}
\usepackage{multirow}
\usepackage{siunitx}
\usepackage[protrusion=true]{microtype}
\usepackage{hyperref}
\usepackage{multicol}
\usepackage{amsmath,amssymb}
\usepackage[a4paper,top=2.5cm,left=2.0cm,right=2.0cm,bottom=2.0cm]{geometry}
\renewcommand{\rmdefault}{ptm}

\sisetup{
  output-decimal-marker = {,}
}

\begin{document}

\title{Evaluación de algoritmos bioinspirados recientes para el Vehicle Routing Problem}
\author{%
  \begin{tabular}{c}
    \textbf{Felipe Ignacio González G.}\\
    Programa de Magíster en Informática\\
    Universidad de Valparaíso, Chile\\
    \href{mailto:felipe.gonzalezggo@postgrado.uv.cl}{felipe.gonzalezggo@postgrado.uv.cl}
  \end{tabular}}
\date{}
\maketitle

{\small
\noindent\textbf{Resumen ---} Se presentó un estudio comparativo riguroso de algoritmos metaheurísticos bioinspirados aplicados al Problema de Ruteo de Vehículos (VRP) sobre instancias Solomon. El análisis incluye dieciséis algoritmos recientes (2016--2025) como WOA, MRFO, FSA y SMA evaluados mediante \textit{benchmarking} masivo con 1.000 ejecuciones por algoritmo. Los resultados establecen cuatro niveles de rendimiento estadísticamente significativos, destacando WOA como superior, seguido por FSA, MRFO y SMA, sin diferencias significativas entre estos cuatro.

\noindent\textbf{Palabras Clave -} metaheurísticas bioinspiradas, \textit{vehicle routing problem}, análisis estadístico, reproducibilidad, diagramas de diferencia crítica.
}

\vspace{1em}
\begin{multicols}{2}

\section{Introducción}
El Problema de Ruteo de Vehículos (VRP) constituye uno de los desafíos fundamentales en investigación operativa y optimización combinatoria. Su naturaleza NP-difícil ha impulsado el desarrollo de múltiples enfoques metaheurísticos, destacando los algoritmos bioinspirados por su capacidad para encontrar soluciones de alta calidad en tiempos razonables.

\section{Metodología}
\subsection{Sistema Experimental}
Implementamos los dieciséis algoritmos en Python 3.11.5, asegurando un marco comparativo justo con interfaces comunes, representación unificada y funciones de evaluación idénticas.

\paragraph{Entorno computacional} Apple MacBook Pro de 16 pulgadas, 2021, equipado con chip Apple M1 Pro (CPU de 10 núcleos y GPU de 16 núcleos) y \SI{16}{\giga\byte} de memoria unificada.

\paragraph{Instancias VRP} Utilizamos seis instancias estándar de la colección Solomon, abarcando distintos tipos de distribución (C: agrupada, R: aleatoria, RC: mixta) y restricciones temporales.

\section{Resultados}
El test de Friedman alineado confirmó diferencias estadísticamente significativas entre los algoritmos evaluados. WOA obtuvo el mejor rendimiento en varias instancias, mientras que SMA demostró la mejor relación coste-calidad.

El análisis estadístico riguroso estableció cuatro grupos de rendimiento con diferencias significativas entre ellos:

\paragraph{Grupo élite (WOA, FSA, MRFO, SMA)} Estos algoritmos demostraron rendimiento superior consistente en todas las instancias, sin diferencias significativas entre ellos.

\section{Conclusiones}
Este estudio presentó una evaluación comparativa rigurosa de dieciséis algoritmos metaheurísticos bioinspirados aplicados al problema VRP. SMA ofrece el mejor equilibrio calidad/coste, requiriendo aproximadamente la mitad del tiempo computacional que WOA con calidad de solución equivalente.

\end{multicols}

\end{document}
